\chapter{存在问题与困难及解决方案}

\section{存在问题与困难}
尽管本课题在多模态蛋白质功能预测方向上已取得初步进展，但在实际研究过程中仍面临以下几个关键问题与挑战：

（1）功能标签的语义表示与分布特性复杂。GO功能标签本身具有多层次、交叉关联、长尾分布等复杂属性，尤其是尾部标签缺乏足够训练样本，容易造成模型预测偏置。尽管本课题引入了大语言文本嵌入模型对标签进行语义编码，但标签嵌入仍存在与蛋白质嵌入对齐不充分的问题，影响语义匹配效果。

(2) 多任务学习的优化冲突。对比学习（强调跨模态对齐）与标签分类学习（强调功能标签判别）的目标存在内在张力，模型在联合训练过程中可能出现优化方向冲突，导致单一任务性能难以同步提升。

(3) 计算资源与实验效率瓶颈。目前的训练机制需要将每个GO功能注释与蛋白质进行一一匹配任务，消耗大量GPU算力与存储资源，限制了超参数搜索范围与消融实验的迭代速度。

\section{解决方案}
针对上述问题与困难，本课题计划采取如下应对策略以优化模型性能、提升研究进度：

(1) 改进功能标签语义建模：在GO语义建模方面，考虑结合图神经网络对GO本体树进行结构化建模，使标签嵌入更好地保留语义层级关系。针对尾部标签，可引入数据增强策略（如语义近邻标签合成）或少样本增强机制，以缓解训练数据稀疏问题。

(2) 提升训练效率与可控性：将训练流程进行分阶段优化，先进行对比学习预训练以学习通用语义对齐能力，再在冻结编码器的基础上进行分类头微调，显著减少参数更新量与训练时间。结合Lora等轻量微调技术，在保持性能的同时降低资源消耗。

总体而言，本课题将持续围绕多模态融合效率、标签语义对齐能力与训练过程可控性等方面展开深入优化研究，力求在提升模型性能的同时，降低训练门槛与资源消耗，进一步增强模型的应用推广价值与科研贡献度。


